% Policy gradient for discounted finite MDPs: a human-readable account of
%   RLTheory/MarkovDecisionProcess/PolicyFamily.lean
%   RLTheory/MarkovDecisionProcess/PolicyGradient.lean
%   RLTheory/Analysis/Calculus/Gradient.lean
%
% Compile with:  pdflatex policy-gradient.tex   (run twice for cross-references)
\documentclass[11pt]{article}

\usepackage[T1]{fontenc}
\usepackage{lmodern}
\usepackage{amsmath,amssymb,amsthm}
\usepackage{mathtools}
\usepackage{booktabs}
\usepackage[margin=1.1in]{geometry}
\usepackage[colorlinks=true,linkcolor=blue,citecolor=blue,urlcolor=blue]{hyperref}

\newtheorem{theorem}{Theorem}[section]
\newtheorem{lemma}[theorem]{Lemma}
\newtheorem{proposition}[theorem]{Proposition}
\newtheorem{corollary}[theorem]{Corollary}
\theoremstyle{definition}
\newtheorem{definition}[theorem]{Definition}
\newtheorem{assumption}[theorem]{Assumption}
\theoremstyle{remark}
\newtheorem{remark}[theorem]{Remark}

% Lean declaration names, printed verbatim (underscores allowed).
\newcommand{\lean}[1]{\texttt{\detokenize{#1}}}
\newcommand{\leanref}[1]{\par{\raggedleft\small Lean:~\lean{#1}\par}}
\emergencystretch=3em

\newcommand{\R}{\mathbb{R}}
\newcommand{\X}{\mathcal{S}}
\newcommand{\U}{\mathcal{A}}
\newcommand{\one}{\mathbf{1}}
\newcommand{\inner}[2]{\langle #1, #2\rangle}
\DeclareMathOperator{\adj}{adj}
\DeclareMathOperator{\sgn}{sgn}

\title{Policy Gradients for Discounted Finite MDPs:\\ A Formalized Account}
\author{Vijaymohan Konda}
\date{November 2026}

\begin{document}
\maketitle

\begin{abstract}
We give a self-contained account of the policy gradient theorem for discounted
Markov decision processes with finite state and action spaces, following the
formalization in Lean~4 contained in the \lean{rl-theory-in-lean} project. The
results are finite-MDP counterparts of the gradient formulas of Chapter~2 of
\emph{Actor-Critic Algorithms} (Konda, 2002), and are the first step towards
formalizing the convergence of actor-critic algorithms. Every definition,
lemma and theorem below has been machine-checked; the name of the corresponding
Lean declaration is given after each statement.
\end{abstract}

\tableofcontents

%======================================================================
\section{Introduction}
%======================================================================

Actor-critic algorithms search for a good policy within a parameterized family
$\{\pi_\theta : \theta \in \R^n\}$ by following an estimate of the gradient of the
overall reward with respect to $\theta$. The starting point of their analysis is
a formula expressing that gradient in terms of quantities that can be estimated
from simulation: the stationary (or discounted) distribution of state-decision
pairs, the state-decision value function $Q_\theta$, and the \emph{score}
$\psi_\theta = \nabla_\theta \log \pi_\theta$.

This note proves that formula for discounted finite MDPs. Two features of the
setting simplify matters considerably compared to the general state spaces of
the thesis.
\begin{enumerate}
\item All objects are finite-dimensional, so value functions are obtained by
  inverting a matrix, and differentiability can be established by elementary
  means (Cramer's rule), with no uniform integrability arguments.
\item We work with \emph{strictly positive} policies. Then the support of the
  induced Markov chain does not depend on $\theta$, so irreducibility and
  aperiodicity need only be checked at a single parameter value.
\end{enumerate}

The proof of the main theorem (Theorem~\ref{thm:pg}) avoids differentiating a
matrix inverse altogether. Instead we differentiate the Bellman equation and
observe that the gradient of the value function solves a Bellman equation of
its own, whose unique solution is given by the same resolvent.

\paragraph{Organization.} Section~\ref{sec:setting} fixes the model.
Section~\ref{sec:policies} introduces parameterized policies and the softmax
family. Section~\ref{sec:values} develops the value function of a fixed policy.
Section~\ref{sec:differentiability} establishes differentiability, and
Section~\ref{sec:gradient} proves the policy gradient theorem.
Section~\ref{sec:lean} lists the correspondence with the Lean code and the
relation to the thesis.

%======================================================================
\section{Setting}\label{sec:setting}
%======================================================================

Throughout, $\X$ and $\U$ are finite nonempty sets of states and decisions
(actions). A vector $x \in \R^{\X}$ is \emph{stochastic} if $x(s) \ge 0$ for all
$s$ and $\sum_s x(s) = 1$; a matrix is \emph{row stochastic} if each of its rows is
stochastic.

\begin{definition}[Finite MDP]\label{def:mdp}
A \emph{finite MDP} consists of
\begin{itemize}
\item transition probabilities $P(s' \mid s, a)$, such that $P(\cdot \mid s, a)$ is
  stochastic for each $(s,a) \in \X \times \U$;
\item a one-stage reward $r : \X \times \U \to \R$;
\item an initial distribution $p_0$ on $\X$;
\item a discount factor $\gamma \in [0, 1)$.
\end{itemize}
\leanref{TabularMDP}
\end{definition}

A (randomized stationary) \emph{policy} is a map $\pi : \X \times \U \to \R$ such
that $\pi(\cdot \mid s)$ is stochastic for every $s$. A policy induces
\begin{align*}
  P_\pi(s, s') &= \sum_{a} \pi(a \mid s)\, P(s' \mid s, a)
    && \text{(state transition matrix)}, \\
  r_\pi(s) &= \sum_{a} \pi(a \mid s)\, r(s, a)
    && \text{(expected one-stage reward)}, \\
  \bar P_\pi\big((s,a), (s',a')\big) &= P(s' \mid s, a)\, \pi(a' \mid s')
    && \text{(state-decision transition matrix)}.
\end{align*}
\leanref{TabularMDP.stateMat, TabularMDP.stateReward, TabularMDP.saMat}

\begin{lemma}\label{lem:rowstoch}
If $\pi$ is a policy, then $P_\pi$ and $\bar P_\pi$ are row stochastic.
\end{lemma}
\begin{proof}
Nonnegativity is clear. For the row sums,
$\sum_{s'} P_\pi(s,s') = \sum_a \pi(a\mid s) \sum_{s'} P(s'\mid s,a) = \sum_a \pi(a\mid s) = 1$,
and similarly
$\sum_{s',a'} \bar P_\pi((s,a),(s',a')) = \sum_{s'} P(s'\mid s,a) \sum_{a'} \pi(a'\mid s') = 1$.
\end{proof}
\leanref{TabularMDP.rowStochastic_stateMat, TabularMDP.rowStochastic_saMat}

\begin{remark}
The existing development in \lean{rl-theory-in-lean} represents an MDP through
probability measures and Markov kernels (\lean{MDPSpec}, \lean{FiniteMDP}). The
real-valued representation above is equivalent for finite spaces and is far more
convenient for calculus. The two are connected: every \lean{MDPSpec} yields a
\lean{TabularMDP} (\lean{MDPSpec.toTabular}), and the state-decision chain of a
\lean{FiniteMDP} used in the existing TD and Q-learning results is exactly
$\bar P_\pi$ (\lean{FiniteMDP.MRP_P_eq_saMat}).
\end{remark}

%======================================================================
\section{Parameterized policies}\label{sec:policies}
%======================================================================

\begin{definition}[Policy family]\label{def:family}
A \emph{policy family} with $n$ parameters is a map
$\theta \mapsto \pi_\theta$ from $\R^n$ to policies such that
\begin{enumerate}
\item $\pi_\theta(a \mid s) > 0$ for all $\theta$, $s$, $a$; and
\item $\theta \mapsto \pi_\theta(a \mid s)$ is differentiable for all $s$, $a$.
\end{enumerate}
The \emph{score} of the family is
\[
  \psi_\theta(s, a) = \nabla_\theta \log \pi_\theta(a \mid s) \in \R^n .
\]
\leanref{PolicyFamily, PolicyFamily.score}
\end{definition}

\begin{lemma}[Likelihood ratio]\label{lem:lr}
For all $\theta$, $s$, $a$,
\[
  \nabla_\theta\, \pi_\theta(a \mid s) = \pi_\theta(a \mid s)\, \psi_\theta(s,a).
\]
\end{lemma}
\begin{proof}
Since $\pi_\theta(a\mid s) > 0$, the chain rule gives
$\psi_\theta(s,a) = \pi_\theta(a\mid s)^{-1}\, \nabla_\theta \pi_\theta(a\mid s)$.
Multiply by $\pi_\theta(a\mid s)$.
\end{proof}
\leanref{PolicyFamily.hasGradientAt_pi}

\begin{lemma}[The score has mean zero]\label{lem:mean-zero}
For all $\theta$ and $s$,
\[
  \sum_{a} \pi_\theta(a \mid s)\, \psi_\theta(s,a) = 0 .
\]
\end{lemma}
\begin{proof}
By Lemma~\ref{lem:lr}, the left side is
$\nabla_\theta \sum_a \pi_\theta(a\mid s) = \nabla_\theta 1 = 0$.
\end{proof}
\leanref{PolicyFamily.sum_smul_score_eq_zero}

\subsection{The softmax family}

Let $\phi : \X \times \U \to \R^n$ be a feature map. The \emph{softmax} (or Gibbs)
policy is
\[
  \pi_\theta(a \mid s) =
  \frac{\exp \inner{\theta}{\phi(s,a)}}{\sum_{b} \exp \inner{\theta}{\phi(s,b)}} .
\]
\leanref{softmax, softmaxFamily}

\begin{proposition}\label{prop:softmax}
The softmax policies form a policy family, and its score is the centered feature:
\[
  \psi_\theta(s,a) = \phi(s,a) - \sum_{b} \pi_\theta(b\mid s)\, \phi(s,b).
\]
\end{proposition}
\begin{proof}
The denominator $Z_\theta(s) = \sum_b \exp\inner{\theta}{\phi(s,b)}$ is positive, so
$\pi_\theta(a\mid s) > 0$; the probabilities clearly sum to one; and
$\pi_\theta(a\mid s)$ is a product of the differentiable functions
$\exp\inner{\theta}{\phi(s,a)}$ and $Z_\theta(s)^{-1}$. For the score, write
$\log \pi_\theta(a\mid s) = \inner{\theta}{\phi(s,a)} - \log Z_\theta(s)$ and note that
$\nabla_\theta Z_\theta(s) = \sum_b \exp\inner{\theta}{\phi(s,b)}\, \phi(s,b)$, so that
$\nabla_\theta \log Z_\theta(s) = \sum_b \pi_\theta(b\mid s)\,\phi(s,b)$.
\end{proof}
\leanref{softmax_pos, softmax_stochastic, differentiable_softmax, softmaxFamily_score}

\subsection{Irreducibility and aperiodicity}

Convergence results for simulation-based algorithms require the Markov chain
of state-decision pairs to be irreducible and aperiodic for every policy in the
family (Assumption~2.4 of the thesis). For strictly positive families it is
enough to check this at a single parameter.

\begin{lemma}\label{lem:pattern}
Let $B$ and $C$ be square matrices with nonnegative entries and the same
positivity pattern, i.e.\ $B_{ij} > 0 \iff C_{ij} > 0$. Then for every $m \ge 0$,
$B^m$ and $C^m$ have the same positivity pattern.
\end{lemma}
\begin{proof}
Induction on $m$, the case $m = 0$ being trivial. Since all entries are
nonnegative,
$(B^{m+1})_{ij} = \sum_k (B^m)_{ik} B_{kj} > 0$ if and only if there is $k$ with
$(B^m)_{ik} > 0$ and $B_{kj} > 0$. By the induction hypothesis and the assumption,
the same condition characterizes $(C^{m+1})_{ij} > 0$.
\end{proof}
\leanref{pow_pos_iff_of_pos_iff}

\begin{proposition}\label{prop:support}
Let $\{\pi_\theta\}$ be a policy family and $\theta_0, \theta \in \R^n$. If
$\bar P_{\pi_{\theta_0}}$ is irreducible (respectively aperiodic), then so is
$\bar P_{\pi_\theta}$.
\end{proposition}
\begin{proof}
Since $\pi_\theta(a'\mid s') > 0$, the entry
$\bar P_{\pi_\theta}((s,a),(s',a')) = P(s'\mid s,a)\,\pi_\theta(a'\mid s')$ is positive
if and only if $P(s'\mid s,a) > 0$, a condition independent of $\theta$. By
Lemma~\ref{lem:pattern}, all powers of $\bar P_{\pi_{\theta_0}}$ and
$\bar P_{\pi_\theta}$ have the same positivity pattern. Irreducibility (for all
$y, z$ some power has a positive $(y,z)$ entry) and aperiodicity (for each $y$ the
set of return times $\{m \ge 1 : (\bar P^m)_{yy} > 0\}$ has greatest common divisor
one) depend only on these patterns.
\end{proof}
\leanref{PolicyFamily.saMat_pow_pos_iff, PolicyFamily.irreducible_saMat, PolicyFamily.aperiodic_saMat}

\begin{remark}
Proposition~\ref{prop:support} does \emph{not} give mixing rates that are uniform
in $\theta$; for the softmax family the chain can mix arbitrarily slowly as
$\|\theta\| \to \infty$. Uniform geometric ergodicity (Assumption~2.5 of the
thesis) therefore remains a genuine assumption for the convergence analysis of
actor-critic algorithms, to be imposed later.
\end{remark}

%======================================================================
\section{Value functions of a fixed policy}\label{sec:values}
%======================================================================

Fix a policy $\pi$ and abbreviate $P = P_\pi$. Define the \emph{resolvent}
\[
  R_\pi = (I - \gamma P_\pi)^{-1},
\]
and the value functions and objective
\begin{align*}
  V_\pi &= R_\pi\, r_\pi, \\
  Q_\pi(s,a) &= r(s,a) + \gamma \sum_{s'} P(s' \mid s,a)\, V_\pi(s'), \\
  J(\pi) &= \sum_s p_0(s)\, V_\pi(s).
\end{align*}
Finally define the normalized discounted state visitation distribution and the
corresponding state-decision occupancy measure
\[
  d_\pi = (1 - \gamma)\, p_0^\top R_\pi, \qquad
  \eta_\pi(s, a) = d_\pi(s)\, \pi(a \mid s).
\]
\leanref{TabularMDP.resolvent, TabularMDP.V, TabularMDP.Q, TabularMDP.J, TabularMDP.visitation, TabularMDP.occupancy}

Since $R_\pi = \sum_{k \ge 0} \gamma^k P_\pi^k$ when the series converges, $d_\pi(s)$ is
the (normalized) expected discounted number of visits to $s$ starting from $p_0$.
We do not use this series; all properties of $R_\pi$ needed below are derived
from two maximum/minimum principles.

\begin{lemma}\label{lem:invertible}
$I - \gamma P_\pi$ is invertible.
\end{lemma}
\begin{proof}
It suffices to show that $(I - \gamma P)v = 0$ implies $v = 0$. Such a $v$ satisfies
$v = \gamma P v$. Choose $s_0$ maximizing $|v(s)|$. Since the row $P(s_0, \cdot)$ is
stochastic,
\[
  |v(s_0)| = \gamma \Big| \sum_{s'} P(s_0,s')\, v(s') \Big|
  \le \gamma \sum_{s'} P(s_0,s')\, |v(s_0)| = \gamma\, |v(s_0)| .
\]
As $\gamma < 1$, this forces $|v(s_0)| = 0$, hence $v = 0$.
\end{proof}
\leanref{TabularMDP.isUnit_one_sub_smul_stateMat}

Consequently $(I - \gamma P)R_\pi = R_\pi(I - \gamma P) = I$, and in particular
\begin{equation}\label{eq:resolvent}
  R_\pi = I + \gamma P_\pi R_\pi .
\end{equation}
\leanref{TabularMDP.one_sub_mul_resolvent, TabularMDP.resolvent_mul_one_sub, TabularMDP.resolvent_apply}

\begin{lemma}[Minimum principle]\label{lem:minprinciple}
If $u, b \in \R^{\X}$ satisfy $u = b + \gamma P_\pi u$ and $b \ge 0$, then $u \ge 0$.
\end{lemma}
\begin{proof}
Choose $s_0$ minimizing $u$. Since $P(s_0,\cdot)$ is stochastic,
$u(s_0) \le \sum_{s'} P(s_0,s')\,u(s')$. Using $b(s_0) \ge 0$ and $\gamma \ge 0$,
\[
  u(s_0) = b(s_0) + \gamma \sum_{s'} P(s_0,s')\, u(s') \ge \gamma\, u(s_0),
\]
so $(1-\gamma)\,u(s_0) \ge 0$ and $u(s_0) \ge 0$. Hence $u \ge u(s_0) \ge 0$.
\end{proof}
\leanref{TabularMDP.nonneg_of_eq_add}

\begin{corollary}\label{cor:resolvent}
The resolvent has nonnegative entries and each of its rows sums to
$(1-\gamma)^{-1}$:
\[
  R_\pi(s,j) \ge 0, \qquad \sum_j R_\pi(s,j) = \frac{1}{1-\gamma} .
\]
\end{corollary}
\begin{proof}
For fixed $j$, the column $u = R_\pi(\cdot, j)$ satisfies $u = e_j + \gamma P u$ by
\eqref{eq:resolvent}, so $u \ge 0$ by Lemma~\ref{lem:minprinciple}. For the row
sums, let $c = (1-\gamma)^{-1}$, so that $c = 1 + \gamma c$. Summing
\eqref{eq:resolvent} over $j$ shows that $w(s) = \sum_j R_\pi(s,j) - c$ satisfies
$w = \gamma P w$, using that $P$ is row stochastic. Applying
Lemma~\ref{lem:minprinciple} with $b = 0$ to both $w$ and $-w$ gives $w = 0$.
\end{proof}
\leanref{TabularMDP.resolvent_nonneg, TabularMDP.sum_resolvent}

\begin{proposition}[Bellman equation]\label{prop:bellman}
$V_\pi = r_\pi + \gamma P_\pi V_\pi$, and
$V_\pi(s) = \sum_a \pi(a\mid s)\, Q_\pi(s,a)$ for every $s$.
\end{proposition}
\begin{proof}
$(I - \gamma P)V_\pi = (I-\gamma P)R_\pi r_\pi = r_\pi$, which rearranges to the
first identity. For the second,
$\sum_a \pi(a\mid s) Q_\pi(s,a) = r_\pi(s) + \gamma \sum_{s'} P_\pi(s,s') V_\pi(s')$
by the definitions of $Q_\pi$ and $P_\pi$, which equals $V_\pi(s)$ by the first
identity.
\end{proof}
\leanref{TabularMDP.bellman, TabularMDP.bellman_apply, TabularMDP.V_eq_sum_Q}

\begin{proposition}\label{prop:visitation}
$d_\pi$ is a probability distribution on $\X$ and $\eta_\pi$ is a probability
distribution on $\X \times \U$.
\end{proposition}
\begin{proof}
Nonnegativity follows from Corollary~\ref{cor:resolvent}. For the total mass,
\[
  \sum_{s'} d_\pi(s') = (1-\gamma) \sum_s p_0(s) \sum_{s'} R_\pi(s,s')
  = (1-\gamma)\cdot\frac{1}{1-\gamma}\sum_s p_0(s) = 1,
\]
and $\sum_{s,a}\eta_\pi(s,a) = \sum_s d_\pi(s)\sum_a \pi(a\mid s) = 1$.
\end{proof}
\leanref{TabularMDP.stochastic_visitation, TabularMDP.stochastic_occupancy}

%======================================================================
\section{Differentiability}\label{sec:differentiability}
%======================================================================

From now on $\{\pi_\theta\}$ is a policy family, and we write
$P_\theta, r_\theta, R_\theta, V_\theta, Q_\theta, J(\theta), d_\theta, \eta_\theta$ for
the objects of Section~\ref{sec:values} associated with $\pi_\theta$.

\begin{lemma}\label{lem:dPdr}
For all $s, s'$,
\[
  \nabla_\theta P_\theta(s,s') = \sum_a P(s'\mid s,a)\, \pi_\theta(a\mid s)\,\psi_\theta(s,a),
  \qquad
  \nabla_\theta r_\theta(s) = \sum_a r(s,a)\, \pi_\theta(a\mid s)\,\psi_\theta(s,a).
\]
\end{lemma}
\begin{proof}
Both are finite sums of $\pi_\theta(a\mid s)$ times constants; apply
Lemma~\ref{lem:lr}.
\end{proof}
\leanref{PolicyFamily.hasGradientAt_stateMat, PolicyFamily.hasGradientAt_stateReward}

\begin{lemma}[Cramer's rule]\label{lem:cramer}
Let $\theta \mapsto B_\theta$ be a matrix-valued map whose entries are
differentiable and with $\det B_\theta \ne 0$ for all $\theta$. Then every entry of
$\theta \mapsto B_\theta^{-1}$ is differentiable.
\end{lemma}
\begin{proof}
By the Leibniz formula
$\det B = \sum_\sigma \sgn(\sigma) \prod_i B_{\sigma(i), i}$, the determinant is a
polynomial in the entries, hence $\theta \mapsto \det B_\theta$ is differentiable.
The entry $(\adj B)_{ij}$ of the adjugate is the determinant of the matrix
obtained from $B$ by replacing its $j$-th row with the $i$-th unit vector; its
entries are entries of $B$ or constants, so it is differentiable as well.
Finally $(B^{-1})_{ij} = (\det B)^{-1} (\adj B)_{ij}$, and $(\det B_\theta)^{-1}$ is
differentiable because $\det B_\theta \ne 0$.
\end{proof}
\leanref{differentiable_det_of_entries, differentiable_inv_of_entries}

\begin{corollary}\label{cor:diffV}
The maps $\theta \mapsto R_\theta(s,j)$, $\theta \mapsto V_\theta(s)$ and
$\theta \mapsto J(\theta)$ are differentiable.
\end{corollary}
\begin{proof}
Apply Lemma~\ref{lem:cramer} to $B_\theta = I - \gamma P_\theta$, whose entries are
differentiable by Lemma~\ref{lem:dPdr} and whose determinant never vanishes by
Lemma~\ref{lem:invertible}. Then $V_\theta(s) = \sum_{s'} R_\theta(s,s')\, r_\theta(s')$ and
$J(\theta) = \sum_s p_0(s) V_\theta(s)$ are finite sums of products of
differentiable functions.
\end{proof}
\leanref{PolicyFamily.differentiable_resolvent, PolicyFamily.differentiable_V, PolicyFamily.differentiable_J}

%======================================================================
\section{The policy gradient theorem}\label{sec:gradient}
%======================================================================

The key observation is that the gradients of the value function satisfy a
Bellman equation whose ``reward'' is $\sum_a \pi_\theta Q_\theta \psi_\theta$, and whose
solution is therefore again given by the resolvent. We first record the
uniqueness statement for vector-valued unknowns.

\begin{lemma}\label{lem:vecbellman}
Let $\pi$ be a policy and let $G, b : \X \to \R^n$ satisfy
$G(s) = b(s) + \gamma \sum_{s'} P_\pi(s,s')\, G(s')$ for all $s$. Then
\[
  G(s) = \sum_{s'} R_\pi(s,s')\, b(s') \qquad \text{for all } s.
\]
\end{lemma}
\begin{proof}
Substituting $b(s') = G(s') - \gamma\sum_k P_\pi(s',k)\,G(k)$ and exchanging the
order of summation,
\begin{align*}
  \sum_{s'} R_\pi(s,s')\, b(s')
  &= \sum_k \Big[ R_\pi(s,k) - \gamma \sum_{s'} R_\pi(s,s') P_\pi(s',k) \Big] G(k) \\
  &= \sum_k \big[R_\pi (I - \gamma P_\pi)\big]_{sk}\, G(k) = G(s). \qedhere
\end{align*}
\end{proof}
\leanref{TabularMDP.eq_resolvent_smul}

\begin{theorem}[Gradient of the value function]\label{thm:gradV}
For every $\theta$ and $s$,
\[
  \nabla_\theta V_\theta(s) = \sum_{s'} R_\theta(s,s') \sum_{a}
    \pi_\theta(a\mid s')\, Q_\theta(s',a)\, \psi_\theta(s',a).
\]
\end{theorem}
\begin{proof}
Write $G(s) = \nabla_\theta V_\theta(s)$, which exists by Corollary~\ref{cor:diffV}.
The Bellman equation (Proposition~\ref{prop:bellman}) holds for every value of
the parameter:
\[
  V_\theta(s) = r_\theta(s) + \gamma \sum_{s'} P_\theta(s,s')\, V_\theta(s').
\]
Differentiating both sides with Lemma~\ref{lem:dPdr} and the product rule,
\begin{align*}
  G(s) &= \sum_a r(s,a)\,\pi_\theta(a\mid s)\,\psi_\theta(s,a)
    + \gamma \sum_{s'} V_\theta(s') \sum_a P(s'\mid s,a)\,\pi_\theta(a\mid s)\,\psi_\theta(s,a) \\
  &\qquad + \gamma \sum_{s'} P_\theta(s,s')\, G(s') \\
  &= \sum_a \pi_\theta(a\mid s)\Big[ r(s,a) + \gamma \sum_{s'} P(s'\mid s,a)\, V_\theta(s')\Big]
    \psi_\theta(s,a) \\
  &\qquad + \gamma \sum_{s'} P_\theta(s,s')\, G(s') \\
  &= b(s) + \gamma \sum_{s'} P_\theta(s,s')\, G(s'),
\end{align*}
where $b(s) = \sum_a \pi_\theta(a\mid s)\, Q_\theta(s,a)\, \psi_\theta(s,a)$. The claim now
follows from Lemma~\ref{lem:vecbellman}.
\end{proof}
\leanref{PolicyFamily.hasGradientAt_V}

\begin{theorem}[Policy gradient theorem, discounted reward]\label{thm:pg}
For every $\theta \in \R^n$,
\[
  \nabla J(\theta) = \frac{1}{1-\gamma} \sum_{s \in \X} \sum_{a \in \U}
    \eta_\theta(s,a)\, Q_\theta(s,a)\, \psi_\theta(s,a).
\]
Equivalently, writing $\inner{f}{g}_\theta = \sum_{s,a}\eta_\theta(s,a) f(s,a) g(s,a)$,
the $i$-th partial derivative is
$\partial_i J(\theta) = (1-\gamma)^{-1}\inner{Q_\theta}{\psi^i_\theta}_\theta$.
\end{theorem}
\begin{proof}
By Theorem~\ref{thm:gradV}, with $b$ as in its proof,
\[
  \nabla J(\theta) = \sum_s p_0(s)\, \nabla_\theta V_\theta(s)
  = \sum_{s'} \Big[\sum_s p_0(s)\, R_\theta(s,s')\Big] b(s')
  = \frac{1}{1-\gamma} \sum_{s'} d_\theta(s')\, b(s'),
\]
and $d_\theta(s')\, b(s') = \sum_a \eta_\theta(s',a)\, Q_\theta(s',a)\, \psi_\theta(s',a)$.
\end{proof}
\leanref{PolicyFamily.hasGradientAt_J}

\begin{corollary}[Baselines]\label{cor:baseline}
For every function $\beta : \X \to \R$,
\[
  \nabla J(\theta) = \frac{1}{1-\gamma} \sum_{s,a}
    \eta_\theta(s,a)\, \big(Q_\theta(s,a) - \beta(s)\big)\, \psi_\theta(s,a).
\]
\end{corollary}
\begin{proof}
The additional term is
$\sum_s d_\theta(s)\,\beta(s) \sum_a \pi_\theta(a\mid s)\,\psi_\theta(s,a)$, which
vanishes by Lemma~\ref{lem:mean-zero}.
\end{proof}
\leanref{PolicyFamily.hasGradientAt_J_baseline}

\begin{remark}
Taking $\beta = V_\theta$ and using Proposition~\ref{prop:bellman} expresses the
gradient through the \emph{advantage} $Q_\theta - V_\theta$. This is the finite-MDP
form of the observation in Chapter~5 of the thesis that the critic need only
learn the projection of $Q_\theta$ onto the span of the score functions, which is
the subject of the next step of the formalization.
\end{remark}

%======================================================================
\section{Correspondence with the Lean code and the thesis}\label{sec:lean}
%======================================================================

\subsection{Lean declarations}

The results above are contained in three files of the
\lean{rl-theory-in-lean} project:
\begin{itemize}
\item \lean{RLTheory/MarkovDecisionProcess/PolicyFamily.lean}: Sections~\ref{sec:setting}--\ref{sec:policies};
\item \lean{RLTheory/MarkovDecisionProcess/PolicyGradient.lean}: Sections~\ref{sec:values}--\ref{sec:gradient};
\item \lean{RLTheory/Analysis/Calculus/Gradient.lean}: sum, product, logarithm
  and exponential rules for gradients, which Mathlib states only for Fr\'echet
  derivatives.
\end{itemize}
All declarations are in the namespace \lean{ReinforcementLearning}. Gradients are
expressed with Mathlib's \lean{HasGradientAt}, on the parameter space
$\R^n$ realized as Mathlib's Euclidean space (abbreviated \lean{E n} in the
project). No declaration depends on any axiom beyond the three standard ones
(\lean{propext}, \lean{Classical.choice}, \lean{Quot.sound}).

\begin{center}
\begin{tabular}{@{}ll@{}}
\toprule
Result & Lean declaration \\
\midrule
Definition~\ref{def:mdp} & \lean{TabularMDP} \\
Definition~\ref{def:family} & \lean{PolicyFamily}, \lean{PolicyFamily.score} \\
Lemma~\ref{lem:lr} & \lean{PolicyFamily.hasGradientAt_pi} \\
Lemma~\ref{lem:mean-zero} & \lean{PolicyFamily.sum_smul_score_eq_zero} \\
Proposition~\ref{prop:softmax} & \lean{softmaxFamily}, \lean{softmaxFamily_score} \\
Proposition~\ref{prop:support} & \lean{PolicyFamily.irreducible_saMat}, \lean{PolicyFamily.aperiodic_saMat} \\
Lemma~\ref{lem:invertible} & \lean{TabularMDP.isUnit_one_sub_smul_stateMat} \\
Lemma~\ref{lem:minprinciple} & \lean{TabularMDP.nonneg_of_eq_add} \\
Corollary~\ref{cor:resolvent} & \lean{TabularMDP.resolvent_nonneg}, \lean{TabularMDP.sum_resolvent} \\
Proposition~\ref{prop:bellman} & \lean{TabularMDP.bellman}, \lean{TabularMDP.V_eq_sum_Q} \\
Proposition~\ref{prop:visitation} & \lean{TabularMDP.stochastic_visitation}, \lean{TabularMDP.stochastic_occupancy} \\
Lemma~\ref{lem:cramer} & \lean{differentiable_inv_of_entries} \\
Lemma~\ref{lem:vecbellman} & \lean{TabularMDP.eq_resolvent_smul} \\
Theorem~\ref{thm:gradV} & \lean{PolicyFamily.hasGradientAt_V} \\
Theorem~\ref{thm:pg} & \lean{PolicyFamily.hasGradientAt_J} \\
Corollary~\ref{cor:baseline} & \lean{PolicyFamily.hasGradientAt_J_baseline} \\
\bottomrule
\end{tabular}
\end{center}

\subsection{Relation to the thesis}

Theorem~\ref{thm:pg} is the finite-MDP counterpart of the discounted-reward
gradient formula of Section~2.4 of the thesis (Theorem~2.13). The differences
are the following.
\begin{itemize}
\item \emph{State and decision spaces.} The thesis allows general (finite,
  countable or continuous) spaces, which requires the bounding functions and
  uniform integrability conditions of Assumptions~2.7--2.9 to justify
  exchanging differentiation and integration. For finite spaces these
  conditions are automatic.
\item \emph{Positivity and smoothness.} Assumption~2.8 of the thesis requires
  $\pi_\theta(a\mid s) > 0$ and continuous differentiability of
  $\theta\mapsto\pi_\theta(a\mid s)$, together with bounds on the score. We make the
  same positivity assumption but require only differentiability; the bounds are
  automatic for finite spaces. The thesis notes that positivity can be weakened
  to ``$\pi_\theta(a\mid s)$ is either positive for all $\theta$ or zero for all
  $\theta$''; that extension is not formalized here. Positivity is also what
  yields Proposition~\ref{prop:support}.
\item \emph{Method of proof.} The thesis differentiates the Poisson (Bellman)
  equation after establishing differentiability of a family of solutions via
  regenerative representations. Here differentiability comes from Cramer's rule
  (Lemma~\ref{lem:cramer}), after which the argument of
  Theorem~\ref{thm:gradV} is the same differentiation of the Bellman equation.
\end{itemize}

\subsection{Next steps}

The remaining steps towards the convergence of actor-critic algorithms are,
in order:
the average-reward counterpart of Theorem~\ref{thm:pg} (Theorem~2.10 of the
thesis), via the fundamental matrix of the chain;
the projection characterization of the gradient and the bias bound of
Lemma~5.3;
an abstract result that $\liminf_k \|\nabla J(\theta_k)\| \le \varepsilon$ for gradient
ascent with vanishing noise and bias at most $\varepsilon$;
and linear stochastic approximation driven by slowly varying Markov chains
(Chapter~3), which together yield the convergence theorem for actor-critic
algorithms (Theorem~5.5).

\begin{thebibliography}{9}
\bibitem{konda2002}
V.~R. Konda.
\newblock \emph{Actor-Critic Algorithms}.
\newblock PhD thesis, Massachusetts Institute of Technology, 2002.
\newblock \url{https://dspace.mit.edu/entities/publication/de41d9e3-f2ba-4327-bdaa-32014dd235ae}.

\bibitem{zhang2026}
S.~Zhang.
\newblock Towards formalizing reinforcement learning theory: a
  Robbins--Siegmund approach.
\newblock \emph{Reinforcement Learning Journal}, 2026.
\newblock \url{https://arxiv.org/abs/2511.03618}.
\end{thebibliography}

\end{document}
